% Part: methods
% Chapter: induction
% Section: introduction

\documentclass[../../../include/open-logic-section]{subfiles}

\begin{document}

\olfileid{mth}{ind}{int}

\olsection{ভূমিকা}

আরোহ একটি গুরুত্বপূর্ণ প্রমাণ-কৌশল। যুক্তিবিদ্যা, তাত্ত্বিক
কম্পিউটারবিজ্ঞান ও গণিতের প্রায় সব ক্ষেত্রেই এর নানা রূপ
ব্যবহৃত হয়। যুক্তিবিদ্যার বহু ফল প্রমাণ করতে আরোহ লাগে।

আরোহকে প্রায়ই অবরোহের বিপরীতে রাখা হয় এবং বিশেষ ঘটনা
থেকে সাধারণ সিদ্ধান্তে পৌঁছানোর অনুমান হিসেবে বর্ণনা করা হয়।
ধরা যাক, আমরা অনেক সবুজ পান্না দেখেছি, কিন্তু পান্না বলব
এমন কোনো অ-সবুজ বস্তু দেখিনি। তখন আমরা সিদ্ধান্ত করতে পারি
যে সব পান্নাই সবুজ। এটি আরোহী অনুমান: অনেক বিশেষ ঘটনা
(এই পান্নাটি সবুজ, ওই পান্নাটি সবুজ, ইত্যাদি) থেকে একটি
সাধারণ দাবি (সব পান্নাই সবুজ) পাওয়া হচ্ছে। \emph{গাণিতিক}
আরোহেও একটি সাধারণ দাবি প্রতিষ্ঠা করা হয়, কিন্তু এই
``সরল আরোহ'' থেকে তার ধরন অনেক আলাদা।

মোটামুটিভাবে, গণিতে আরোহ-প্রমাণ দেখায় যে এক ধরনের সব
গাণিতিক বস্তুর একটি নির্দিষ্ট ধর্ম আছে। সবচেয়ে সহজ ক্ষেত্রে
সেই বস্তুগুলি স্বাভাবিক সংখ্যা। তখন আরোহ-প্রমাণ দিয়ে
দেখানো হয় যে সব স্বাভাবিক সংখ্যার ধর্মটি আছে। তার জন্য
দুটি বিষয় দেখাতে হয়:
\begin{enumerate}
    \item $0$-এর ধর্মটি আছে, এবং
    \item যখনই কোনো সংখ্যা~$k$-এর ধর্মটি থাকে, তখন~$k+1$-এরও থাকে।
\end{enumerate}
স্বাভাবিক সংখ্যার উপর আরোহ ব্যবহার করে প্রায়ই এমন
গাণিতিক বস্তু সম্পর্কেও সাধারণ দাবি প্রমাণ করা যায়,
যাদের সঙ্গে সংখ্যা নির্ধারিতভাবে যুক্ত করা যায়। যেমন,
প্রতিটি সসীম সেটের সসীমসংখ্যক~$n$ উপাদান আছে। আরোহ দিয়ে
যদি দেখাতে পারি, প্রতিটি সংখ্যা~$n$-এর জন্য ``$n$-আকারের
সব সসীম সেট \dots'' ধর্মটি সত্য, তবে সব সসীম সেট সম্পর্কে
একটি ফলই প্রমাণ করা হবে।

\emph{আরোহী সংজ্ঞায়িত} গাণিতিক বস্তুর ক্ষেত্রেও আরোহকে
সাধারণীকরণ করা যায়। যেমন, প্রথম-ক্রমের যুক্তিবিদ্যার
অভিব্যক্তির মতো আনুষ্ঠানিক ভাষার অভিব্যক্তিগুলি আরোহীভাবে
সংজ্ঞায়িত হয়। \emph{গঠনগত আরোহ} দিয়ে এমন সব অভিব্যক্তি
সম্পর্কে ফল প্রমাণ করা যায়। বিশেষ করে যুক্তিবিদ্যায়
গঠনগত আরোহ অত্যন্ত উপযোগী এবং বহুল ব্যবহৃত।

\end{document}
